% =====================================================================
%  CARNET D'ENTRAINEMENT — FICHE 6 — THÈME 5
%  Niveau MOYEN — Exercices
% =====================================================================
\documentclass[11pt,a4paper]{article}
\usepackage[utf8]{inputenc}
\usepackage[T1]{fontenc}
\usepackage{lmodern}
\usepackage[french]{babel}
\usepackage{amsmath,amssymb,mathtools}
\usepackage{icomma}
\usepackage[version=4]{mhchem}
\usepackage{siunitx}
\sisetup{output-decimal-marker={,}}
\usepackage{xcolor}
\usepackage{colortbl}
\usepackage{tikz}
\usepackage[most]{tcolorbox}
\usepackage{enumitem}
\usepackage{array}
\usepackage{booktabs}
\usepackage{fancyhdr}
\usepackage[hidelinks]{hyperref}
\usetikzlibrary{arrows.meta,positioning,calc,shapes.geometric}
\usepackage[a4paper,margin=2.1cm,headheight=26pt,headsep=8pt]{geometry}

\definecolor{primary}{RGB}{146,95,160}
\definecolor{primarydark}{RGB}{92,52,104}
\definecolor{excol}{RGB}{0,135,90}
\definecolor{atomO}{RGB}{220,55,45}
\definecolor{bondcol}{RGB}{60,60,70}
\definecolor{piecol}{RGB}{198,93,9}

\tcbset{boxbase/.style={enhanced,breakable,boxrule=0.9pt,arc=2.5pt,
   left=3mm,right=3mm,top=2.2mm,bottom=2.2mm,
   fonttitle=\bfseries,coltitle=white,toptitle=0.6mm,bottomtitle=0.6mm}}
\newtcolorbox[auto counter]{exo}[1][]{boxbase,colback=primary!5,colframe=primary,
   title={\textsc{Exercice}~\thetcbcounter\ifx\relax#1\relax\relax\else\textmd{ —\itshape\ #1}\fi}}

\pagestyle{fancy}\fancyhf{}
\fancyhead[L]{\small\textcolor{primarydark}{\textbf{Fiche 6 · Thème 5 — La résonance}}}
\fancyhead[R]{\small\textcolor{primarydark}{\textsc{Moyen}}}
\fancyfoot[C]{\small\thepage}
\renewcommand{\headrulewidth}{0.8pt}
\renewcommand{\headrule}{\hbox to\headwidth{\color{primary}\leaders\hrule height \headrulewidth\hfill}}
\setlength{\parindent}{0pt}\setlength{\parskip}{3pt}

\begin{document}
\begin{center}
{\Large\bfseries\color{primarydark}Thème 5 — Niveau Moyen}\\[2pt]
{\small\itshape Justifier chaque réponse. La double flèche de résonance se note $\longleftrightarrow$.}
\end{center}
\medskip

\begin{exo}[l'ozone \ce{O3}]
La molécule d'ozone est coudée, avec un atome d'oxygène central lié à deux oxygènes.
\begin{enumerate}[label=\alph*),leftmargin=2em,itemsep=1pt]
  \item Écrire les \textbf{deux} formes limites de l'ozone reliées par la flèche $\longleftrightarrow$.
  \item Expliquer pourquoi les deux liaisons \ce{O-O} \emph{réelles} sont identiques, et préciser
        comment se compare leur longueur à celle d'une liaison simple et d'une liaison double.
\end{enumerate}
\end{exo}

\begin{exo}[l'ion carbonate \ce{CO3^2-}]
\begin{enumerate}[label=\alph*),leftmargin=2em,itemsep=1pt]
  \item Expliquer pourquoi il existe \textbf{trois} formes limites équivalentes pour cet ion.
  \item En déduire ce que l'on peut dire des trois liaisons \ce{C-O} réelles.
  \item La charge globale de l'ion est $-2$. Comment se répartit-elle en moyenne sur les trois
        oxygènes ? Donner la charge portée par chaque oxygène.
\end{enumerate}
\end{exo}

\begin{exo}[l'ion nitrate \ce{NO3-}]
L'ion nitrate a la même structure « à trois oxygènes autour d'un atome central » que le carbonate.
\begin{enumerate}[label=\alph*),leftmargin=2em,itemsep=1pt]
  \item Par analogie avec \ce{CO3^2-}, indiquer le nombre de formes limites équivalentes de \ce{NO3-}.
  \item En déduire ce que l'on prévoit pour les longueurs des trois liaisons \ce{N-O}.
\end{enumerate}
\end{exo}

\begin{exo}[deux flèches à ne pas confondre]
On considère d'une part la résonance de l'ozone \ce{O=O-O <-> O-O=O}, d'autre part l'équilibre
d'ionisation de l'eau \ce{2 H2O <=> H3O+ + OH-}.
\begin{enumerate}[label=\alph*),leftmargin=2em,itemsep=1pt]
  \item Lequel de ces deux schémas met en jeu \emph{une seule} espèce ? Lequel en met en jeu
        \emph{deux} (ou plus) ?
  \item Expliquer la différence de \textbf{sens physique} entre les symboles $\longleftrightarrow$
        et $\rightleftharpoons$.
\end{enumerate}
\end{exo}

\begin{exo}[l'ion perchlorate \ce{ClO4-}]
Dans \ce{ClO4-}, les quatre oxygènes sont expérimentalement \emph{équivalents} et les quatre liaisons
\ce{Cl-O} ont la même longueur.
\begin{enumerate}[label=\alph*),leftmargin=2em,itemsep=1pt]
  \item La « meilleure » structure de Lewis place la charge $-1$ sur \emph{un} oxygène précis et des
        doubles liaisons sur les autres. En quoi cela semble-t-il \emph{contredire} la réalité ?
  \item Expliquer comment la résonance résout ce paradoxe, et donner le nombre de formes limites
        équivalentes.
\end{enumerate}
\end{exo}

\end{document}
